\begin{textsample}{NEG Dim 5 – global\_north\_2023\_12 – Score -1.00 – global\_north\_2023\_12/104448743.txt}  \label{ex:f5_neg_053}
Hamas and Israeli officials appear open to discussing a \textbf{weeklong} truce that could free as many as 40 hostages -- including women, the elderly and those who are sick.
All three major Israeli TV networks presented the Jewish state ' s negotiating goals on Tuesday, revealing that their government wants to see Hamas free specific hostages, including women, seniors, and those with physical or mental illnesses.
In exchange for the hostages, Israel is allegedly prepared to engage in a \textbf{weeklong} cease-fire in Gaza and the release of Palestinian prisoners in its jails, the Times of Israel reports.
The deal would also seek to expand the humanitarian zones in the Palestinian enclave, as well as the amount of aid being delivered to Gaza.
It is similar to the deal that saw a week of peace between the warring armies in late November, which ultimately fell apart after Hamas and Israel could not come to an agreement on an eighth wave of hostage exchanges.
Hamas had released 105 people during a \textbf{weeklong} cease-fire in late November.
Hamas Press Service/UPI/Shutterstock President @ @ @ @ @ @ @ @ @ @ Israel and Hamas to reach a deal that could see more of the estimated 128 hostages freed."We ' re pushing it.
There ' s no expectation at this point, but we are pushing,"Biden said as he described reports of the Palestinian death toll likely surpassing 20,000 later today as"tragic."The reports of the Jewish state ' s openness to the deal came as Israeli President Isaac Herzog claimed his nation was prepared for a second pause in fighting -- as long as Hamas adheres to the agreed-upon rules -- during a meeting with ambassadors from 80 countries on Tuesday.
Israel, which has decimated Hamas terror infrastructure, said the fighting will not end until the group is destroyed.
AFP via Getty Images A source close to Hamas has since confirmed an Axios report that the group is open to the deal, with its top leader, Ismail Haniyeh, traveling to Egypt on Wednesday.
Hamas said in a statement that Haniyeh would be discussing the war with Egyptian officials, but @ @ @ @ @ @ @ @ @ @ negotiators were the ones who brokered the first truce, and have since been attempting to restore peace in Gaza after intense fighting resumed on Dec. 1.
Despite the signs of another temporary truce on the horizon, high tension remains between the two sides, with Hamas signaling that a deal will not happen until Israel leaves Gaza, and the Jewish state reiterating that the fighting will not end until the terrorist group is destroyed.

% matched lemmas: weeklong
\end{textsample}
